\begin{enumerate}

\item La figura establece qué partes pertenecen a la sección, la dirección de y y z y las dimensiones en milímetros.

\item El área total suma las áreas macizas y resta los huecos; los primeros momentos proporcionan las coordenadas del centroide.

\item Calcular I\_z por suma de las inercias propias de las piezas más A\_i por el cuadrado de cada distancia a G (Steiner).

\item El verificador geométrico de Studium contrasta las inercias por otra vía: integración directa sobre los contornos.

\item Si los ejes no son principales, obtener los valores propios de la matriz de inercias sin suponer que I\_yz = 0.

\end{enumerate}

\noindent\textbf{Operaciones y resultados intermedios.}

\begin{enumerate}

\item Área algebraica de la sección, descontando los huecos.

\[\fitmath{y_{1}=\left(10-0\right)\cdot \left(100-0\right)+\left(60-10\right)\cdot \left(100-90\right)\approx 1500\,\mathrm{mm^{2}}}\]

\item Centroide en el eje vertical mediante primeros momentos.

\[\fitmath{y_{2}=\frac{\frac{\left(10-0\right)\cdot \left(100-0\right)\cdot \left(0+100\right)}{2}+\frac{\left(60-10\right)\cdot \left(100-90\right)\cdot \left(90+100\right)}{2}}{\left(10-0\right)\cdot \left(100-0\right)+\left(60-10\right)\cdot \left(100-90\right)}\approx 65\,\mathrm{mm}}\]

\item Momento de inercia I\_z mediante el teorema de Steiner.

\[\fitmath{y_{3}=\frac{\left(10-0\right)\cdot {\left(100-0\right)}^{3}}{12}+\left(10-0\right)\cdot \left(100-0\right)\cdot {\left(\frac{0+100}{2}-65\right)}^{2}+\frac{\left(60-10\right)\cdot {\left(100-90\right)}^{3}}{12}+\left(60-10\right)\cdot \left(100-90\right)\cdot {\left(\frac{90+100}{2}-65\right)}^{2}\approx 1512500\,\mathrm{mm^{4}}}\]

\item Momento principal máximo de inercia.

\[\fitmath{y_{4}=50000\cdot \sqrt{202}+962500\approx 1673134\,\mathrm{mm^{4}}}\]

\end{enumerate}
